\section{Introduction}\label{sec:intro}

Commercial and institutional buildings in many markets pay a demand charge on their highest metered power in a billing month in addition to an energy charge per kilowatt-hour. A behind-the-meter battery can lower that charge by discharging during the peak, and it can also shift energy from expensive to cheap hours~\cite{oudalov2007sizing,leadbetter2012battery,hanna2014dispatch,papadopoulos2020peak}. Both uses need a forecast of the building's load, and the forecast is usually chosen the way forecasts are chosen everywhere: by its mean absolute error (MAE), its root mean squared error or a percentage error on held-out data~\cite{hong2016probabilistic,wang2019smartmeter,hyndman2006another}. Large benchmarks of building load forecasting rank models in exactly this way~\cite{emami2023buildingsbench,miller2020gepiii}.

There are reasons to doubt that accuracy ranks forecasts by what they are worth to a battery. A demand charge depends on a single hour per month, and an error that makes the forecast too low at a charging hour can set a new peak, while an error of the same size in the other direction only keeps charge in reserve. That the value of a forecast depends on the user's loss, not only on its error, is an old idea in forecast verification~\cite{murphy1993good}; a point forecast is optimal only relative to a scoring function~\cite{gneiting2011point,gneiting2007scoring}, and the value of probabilistic forecasts for decisions under asymmetric penalties is well documented in energy trading~\cite{pinson2007trading}. Decision-focused learning trains forecasters on the cost of the decision they feed instead of on their error~\cite{donti2017task,elmachtoub2022spo,mandi2024dfl,sadana2025survey}, and recent work applies the idea to demand-charge storage~\cite{kim2025neural} and to photovoltaic-battery scheduling in twenty buildings~\cite{depoortere2026dfl}. What is missing is a measurement, on many buildings, of how far the usual accuracy metric is from the decision value of the forecasts that practitioners actually use, of which properties of the error drive the gap, and of whether a simple asymmetric score closes it.

This paper provides that measurement. We schedule a battery day-ahead in each of \res{nKept} buildings of the Building Data Genome 2 data set~\cite{miller2020bdg2} from five standard forecasters, execute every schedule against the actual 2017 load, and compare the share of the perfect-information saving each forecaster captures with its accuracy. The analysis plan, with its primary hypothesis and test, was committed to the study's version-controlled repository before the data were downloaded; three amendments, the last made after the results and an independent review and including a bug fix, are disclosed, and the registered test is reported on the registered simulation. Synthetic forecasts with exactly the MAE of the best model but a single type of error each (uniform over- and under-forecasts, independent and correlated noise, a late ramp) separate which errors matter. Each schedule is executed in two ways, open loop and with a real-time guard that stops the battery from charging into the month's peak, which separates the value of the forecast from that of the controller.

Under the demand charge, the pre-specified test found no agreement between the ranking of forecasters by MAE and their ranking by saving: the mean Kendall rank correlation across buildings was \res{voneMean} (95\% CI \res{voneLo} to \res{voneHi}). The sign of the error explains the gap. The 0.9-quantile forecaster, the least accurate on average, captured the largest saving; at equal MAE a uniform over-forecast kept \res{shT1Sscale} of the attainable saving while a uniform under-forecast lost money when executed open loop; and ranking by the 0.9 pinball loss agreed with value ($\tau=\res{pinOpenMean}$), as did MAE among the four unbiased forecasters ($\tau=\res{fourOpenMean}$). The controller mattered more than the forecaster: open-loop schedules raised the bill in \res{negOpenFourLo}--\res{negOpenFourHi}\% of buildings, whereas the guarded schedules saved money and beat a forecast-free rule in \res{beatReactLo}--\res{beatReactHi}\% of buildings. The saving at stake is small, \res{pctT1PI}\% of the median bill with perfect information, which is why errors of a few per cent in the wrong direction can erase it. Without a demand charge every forecaster captured the full saving. Section~\ref{sec:related} reviews related work, Section~\ref{sec:model} the data and decision model, Section~\ref{sec:design} the forecasters and pre-specified design, Section~\ref{sec:results} the results, and Section~\ref{sec:discussion} their implications and limits.
